% content/c00-theorie/de.tex — Kapitel 0: Was eine physikalische Theorie liefert
% (deutsche Leitfassung)
%
% Migration der ersten Vorlesung latex/lectures/de/vl01.tex (Vorspann, Abschnitt
% „Der Geburt der Quantenmechanik“, Unterabschnitt „Vorgeschichte“ und der graue Kasten
% „Jedes Experiment in drei Schritten“).  Wohin jede Gleichung, jeder Kasten und jede
% Abbildung von vl01 gegangen ist, steht in evidence/migrations/maps/vl01-map.json.
%
% Rahmung: docs/plan/BRIEF.md, Abschnitt „The physics approach“ (Zeilen 38, 40, 44, 45,
%   50, 90).  Sie ist TJOs Festlegung (author-stipulated) und wird hier nur zitiert und
%   angeordnet.  Die deutschen Fassungen der zitierten Sätze sind Übersetzungen des
%   Kapitelautors; jede ist mit „GAP-TJO“ markiert.
% Physik der Vorgeschichte: einzige Quelle ist TJOs Skript (Quellenschlüssel
%   tpc-script-v0.2); jede Formel daraus ist source-backed mit dem Geltungsbereich „wie im
%   Vorlesungsskript des Autors angegeben; keine externe Quelle geprüft“.  Zahlenwerte
%   der Konstanten: CODATA 2022.
% Herleitungen (Grenzform, Schranken, Schwelle): evidence/derivations/c01/DERIVATION.md,
%   Teil „Kapitel 0“.
%
% Schreibweise: Das Skript schreibt U(nu) für die spektrale Energiedichte; hier steht
% u(nu, T), weil U in der Registry mehrfach belegt ist (COL-U) und die Temperatur als
% Argument gebraucht wird.

\chapter{Was eine physikalische Theorie liefert}\label{c00}

Dieses Kapitel sagt zwei Dinge. Erstens: was in diesem Kurs unter einer physikalischen
Theorie verstanden wird und in welcher Reihenfolge der Kurs seine Theorien behandelt.
Zweitens: an welchen Befunden die klassische Physik gescheitert ist, sodass eine neue
Theorie nötig wurde. Gerechnet wird hier wenig. Die Mathematik beginnt in Kapitel~1 mit
der Wahrscheinlichkeitsrechnung.

\section{Die Rahmung des Kurses}\label{c00.s1}

Der Kurs folgt einer Rahmung, die sein Autor festgelegt hat. Sie wird nicht hergeleitet;
sie sagt, wie der Stoff angeordnet und in welcher Sprache er formuliert wird. Ihre ersten
Sätze lauten:
% GAP-TJO: Die deutsche Fassung der drei folgenden Sätze ist eine Übersetzung aus
% docs/plan/BRIEF.md:44, 45 und 50 und nicht von TJO geschrieben; der englische Wortlaut
% steht wörtlich in der englischen Ausgabe und in evidence/claims.json (G-0022).
\begin{equation}\label{c00.eq.liefert}\claim{c00.cl.liefert}{author-stipulated}
  \begin{aligned}
    &\text{\tpcifenglish{A physical theory supplies data for}{Eine physikalische Theorie liefert Daten für}} \\
    &\text{\tpcifenglish{preparations + dynamics + measurement.}{Präparationen + Dynamik + Messung.}}
  \end{aligned}
\end{equation}
\begin{equation}\label{c00.eq.probabilistisch}\claim{c00.cl.probabilistisch}{author-stipulated}
  \begin{aligned}
    &\text{\tpcifenglish{All physical theories are}{Alle physikalischen Theorien sind}} \\
    &\text{\tpcifenglish{probabilistic theories (I think of deterministic}{probabilistische Theorien (deterministische fasse ich}} \\
    &\text{\tpcifenglish{as a subset of probabilistic).}{als Teilmenge der probabilistischen auf).}}
  \end{aligned}
\end{equation}
\begin{equation}\label{c00.eq.p0}\claim{c00.cl.p0}{author-stipulated}
  \begin{aligned}
    &\text{\tpcifenglish{Physical theories are probabilistic theories:}{Physikalische Theorien sind probabilistische Theorien:}} \\
    &\text{\tpcifenglish{predictions in the form of probabilities.}{Vorhersagen in Form von Wahrscheinlichkeiten.}}
  \end{aligned}
\end{equation}
Der Brief nennt Satz~\eqref{c00.eq.p0} P0. Der Kapitelautor liest ihn als Auftakt der
schrittweisen Spezialisierung der Postulate; diese verbindende Formulierung steht zur
Entscheidung des Autors aus (G-0023).

Satz~\eqref{c00.eq.liefert} nennt drei Zutaten. Der nächste Abschnitt zeigt, wo sie in
einem Experiment vorkommen. Satz~\eqref{c00.eq.probabilistisch} sagt, dass deterministische
Theorien als Teil der probabilistischen aufgefasst werden. Kapitel~1 beweist eine
Proposition über Punktgewichte auf endlichen Ergebnismengen. Ihre Deutung als Beispiel
dieser Festlegung ist ein Vorschlag des Kapitelautors und wartet auf die Entscheidung
des Autors (G-0023); eine Definition deterministischer Theorien samt Dynamik steht aus.

\section{Jedes Experiment in drei Schritten}\label{c00.s2}

Die Vorlesung, aus der dieser Kurs entstanden ist, beschreibt Experimente mit einem
einzigen Schema \srcref{tpc-script-v0.2}{vl01, Kasten „Jedes Experiment in drei Schritten“}:
\begin{equation}\label{c00.eq.drei-stadien}\claim{c00.cl.drei-stadien}{source-quoted}
  \begin{aligned}
    &\text{\tpcifenglish{To the theorist all experiments in physics}{Für den Theoretiker sehen alle Experimente}} \\
    &\text{\tpcifenglish{look the same; they consist of three stages.}{der Physik gleich aus; sie bestehen aus drei Stadien.}}
  \end{aligned}
\end{equation}
\begin{enumerate}
  \item \emph{Präparation.} Etwas wird hergerichtet: ein Würfel, im Labor etwa ein
    Lichtstrahl.
  \item \emph{Durchführung.} Das Experiment läuft.
  \item \emph{Messung.} Am Ende schauen wir hin und erhalten eines der möglichen
    Ergebnisse.
\end{enumerate}
Abbildung~\ref{c00.fig.experiment} zeigt das Schema. Ein Experimentalphysiker würde diese
Vereinfachung nicht anerkennen; abstrakt sieht aber jedes Experiment so aus.

\begin{tpcfigure}{c00.fig.experiment}
  \caption{Die drei Stadien eines Experiments. Die Klammer fasst die ersten beiden zu
    einer einzigen Präparation zusammen. Die Präparation ist durchgezogen gezeichnet, die
    Durchführung mit dickerem Strich, die Messung gestrichelt.}
\end{tpcfigure}

Die ersten beiden Schritte lassen sich zu \emph{einer} komplizierteren Präparation
zusammenfassen. Ein gezinkter Würfel ist nichts anderes als ein normaler Würfel, an dem man
vorher etwas gemacht hat. Deshalb kommt man oft mit zwei Begriffen aus: Präparation und
Messung. Der Kapitelautor stellt diese Stadien den drei Zutaten in
Satz~\eqref{c00.eq.liefert} gegenüber: Präparation, Dynamik, Messung. Diese Zuordnung
wartet auf die Entscheidung des Autors (G-0023).
% GAP-TJO: Dass die „Durchführung“ des Skripts (vl01.tex:122) die „dynamics“ des Briefs
% (BRIEF.md:44) ist, steht in keiner der beiden Quellen ausdrücklich; die Gegenüberstellung
% ist die des Kapitelautors (G-0023).

Die folgende Definition ist aus Plávala übernommen; die Abweichungen des Kursautors
von Ludwig-Werner sind noch nicht formuliert (G-0023).

\tpcdef{D1}

Die folgende Erläuterung ist die Lesart des Kapitelautors und wartet auf die Entscheidung
des Autors (G-0023). Eine Präparation ist zunächst etwas, das man \emph{tut}: eine Liste
von Anweisungen. Für die Vorhersagen zählt, welche Wahrscheinlichkeiten die anschließenden
Ja-Nein-Messungen liefern. Zwei Verfahren, die auf alle Ja-Nein-Fragen dieselben
Wahrscheinlichkeiten liefern, beschreibt die Theorie durch denselben Zustand. Die Ja-Nein-Fragen, von
denen die Definition spricht, sind die einfachsten Messungen: solche mit nur zwei möglichen
Ergebnissen. In Kapitel~1 sind es die Indikatorfunktionen, und Kapitel~2 definiert sie
allgemein.

\section{Vom Kleinen zum Großen}\label{c00.s3}

Auch die Reihenfolge des Stoffs ist festgelegt:
% GAP-TJO: Deutsche Fassung der beiden folgenden Sätze (BRIEF.md:38 und 90) und der
% Beschriftungen in Abbildung c00.fig.leiter (BRIEF.md:40): Übersetzung des Kapitelautors
% (G-0022).  Im Brief ist das Wort „effective“ hervorgehoben.
\begin{equation}\label{c00.eq.leiter}\claim{c00.cl.leiter}{author-stipulated}
  \begin{aligned}
    &\text{\tpcifenglish{Our entire presentation works like this:}{Unsere ganze Darstellung geht so vor:}} \\
    &\text{\tpcifenglish{we go from microscopic to macroscopic.}{Wir gehen vom Mikroskopischen zum Makroskopischen.}}
  \end{aligned}
\end{equation}
Abbildung~\ref{c00.fig.leiter} zeigt die Stufen mit ausgewählten Auszügen aus der
Festlegung des Autors; die deutschen Beschriftungen sind Übersetzungen des Kapitelautors
und warten auf den Wortlaut des Autors (G-0022). Von Stufe zu
Stufe wird „herausgezoomt“; an einer Stelle geht dabei ausdrücklich Information verloren.
% GAP-TJO: Welche Operation hinter jedem Pfeil der Leiter steht (PLAN.md 13.1, A11:
% Modellwahl, Erweiterung, Vergröberung mit Grenzübergang), ist hier nicht ausgeführt:
% Der Brief nennt sie nicht, und die Kapitel, die die Schritte gehen, sind noch nicht
% geschrieben (G-0023).

\begin{tpcfigure}{c00.fig.leiter}
  \caption{Die Leiter der Theorien dieses Kurses, von unten (mikroskopisch) nach oben
    (makroskopisch). Die Beschriftungen der sechs Stufen und der fünf Pfeile sind
    ausgewählte Auszüge aus der Festlegung des Autors, auf die sich
    Satz~\eqref{c00.eq.leiter} bezieht. Die letzte Stufe nennt der Brief vollständig
    „Thermodynamik als die ultimative effektive Theorie der gesamten Physik auf
    makroskopischer Skala“; auch diese Übersetzung wartet auf den Wortlaut des Autors
    (G-0022). GPT
    steht für „general probabilistic theory“, QM für Quantenmechanik. Der Kurs beginnt
    unten: Kapitel~1 und~2 behandeln Stufe~1.}
\end{tpcfigure}

Zu jeder Stufe gehört derselbe Gedanke:
\begin{equation}\label{c00.eq.effektiv}\claim{c00.cl.effektiv}{author-stipulated}
  \begin{aligned}
    &\text{\tpcifenglish{At each stage the student learns that}{Auf jeder Stufe lernt man, dass}} \\
    &\text{\tpcifenglish{every theory is an effective theory of the one below it:}{jede Theorie eine effektive Theorie der darunterliegenden ist:}} \\
    &\text{\tpcifenglish{an approximation to the fundamental underlying physics.}{eine Näherung an die grundlegende Physik darunter.}}
  \end{aligned}
\end{equation}
Was „herauszoomen“ bei einem bestimmten Schritt bedeutet, ist Gegenstand des Kapitels,
das diesen Schritt geht. Hier steht nur die Reihenfolge.

\section{Warum eine neue Theorie nötig war}\label{c00.s4}

Die Vorlesung beginnt mit einer Feststellung: Wir befinden uns in einem Quanten-Universum.
Bis 1900 blieben dessen Quantenaspekte unbemerkt, weil die klassische Physik eine sehr gute
Annäherung an die Physik des Alltags liefert. Erklärungsbedürftig ist deshalb nicht die
Quantenmechanik, sondern die Güte der klassischen Näherung. Das Stichwort dazu heißt
\emph{Dekohärenz}; sie gehört in die fortgeschrittene Quantenmechanik und wird hier nicht
behandelt \srcref{tpc-script-v0.2}{vl01, Abschnitt 1}. Der Kurs arbeitet mit der Annahme
\begin{equation}\label{c00.eq.annahme}\claim{c00.cl.annahme}{source-quoted}
  \begin{aligned}
    &\text{\tpcifenglish{Our assumption is that quantum mechanics}{Unsere Annahme ist, dass die Quantenmechanik}} \\
    &\text{\tpcifenglish{is the correct description of the universe.}{die korrekte Beschreibung des Universums ist.}}
  \end{aligned}
\end{equation}

Zur Methode: Man kann Quantenmechanik so lehren, dass man erst die falschen, klassischen
oder halbklassischen Modelle durchrechnet und dann zeigt, dass sie nicht funktionieren.
Diesen Weg geht der Kurs nicht. Das Bohrsche Atommodell wird nirgends durchgerechnet. Die
Befunde, an denen die klassische Physik gescheitert ist, werden trotzdem genannt: zwei
davon mit ihrer Formel, die übrigen als Liste. Keiner wird in diesem Kapitel hergeleitet.
Alle Aussagen dieses Abschnitts über Experimente stammen aus dem Vorlesungsskript und sind
hier nicht gegen eine weitere Quelle geprüft.

\subsection{Schwarzkörperstrahlung}\label{c00.s4.1}

Das Skript beschreibt einen Schwarzkörper als Körper, der keine einfallende Strahlung
reflektiert und bei einer Temperatur $\Temp$ wegen der thermisch angeregten Atome dennoch
strahlt. Als übliche Realisierung nennt es eine Kiste mit einem kleinen Loch, aus dem
Schwarzkörperstrahlung austritt (Abbildung~\ref{c00.fig.hohlraum}). Die Beschreibung ist
eine Aussage der Vorlesung; die physikalische Begründung der Hohlraumidealisierung
wartet auf eine Quelle oder Formulierung des Autors (G-0024).

\begin{tpcfigure}{c00.fig.hohlraum}
  \caption{Die im Skript genannte Realisierung eines Schwarzkörpers: eine Kiste mit einem
    kleinen Loch. Die gezeichneten Reflexionen treffen das Loch nicht; gezeigt sind die
    ersten sechs Reflexionen eines Strahls nach dem Reflexionsgesetz. Das Skript bezeichnet
    die aus dem Loch austretende Strahlung als Schwarzkörperstrahlung. Die Zeichnung zeigt
    den Strahlengang; die physikalische Begründung der Idealisierung steht aus (G-0024).}
\end{tpcfigure}

Das Spektrum dieser Strahlung ist mit klassischer Physik nicht erklärbar. Das richtige
Spektrum liefert die Annahme, dass Licht der Frequenz $\nu$ aus Lichtquanten der Energie
\begin{equation}\label{c00.eq.lichtquant}\claim{c00.cl.lichtquant}{source-backed}
  E = \hPlanck \nu
\end{equation}
besteht. Die Konstante $\hPlanck$ heißt Plancksches Wirkungsquantum. Das Strahlungsgesetz
selbst gibt an, wie viel Energie pro Volumen auf die Frequenzen zwischen $\nu$ und
$\nu + \dd\nu$ entfällt:
\begin{equation}\label{c00.eq.planck}\claim{c00.cl.planck}{source-backed}
  u(\nu, \Temp)\,\dd\nu
  = \frac{8\pi \hPlanck \nu^3}{\clight^3}\,\frac{\dd\nu}{e^{\hPlanck\nu/(\kB \Temp)} - 1} .
\end{equation}
Darin ist $\clight$ die Lichtgeschwindigkeit und $\kB$ die Boltzmann-Konstante. Der Kurs
leitet dieses Gesetz im Teil über statistische Physik selbst her, aus der Annahme, dass
Licht quantisiert ist. Hier wird es nur benutzt.

In der theoretischen Physik verwendet man statt $\hPlanck$ meist die \emph{reduzierte
Planck-Konstante}
\begin{equation}\label{c00.eq.hbar}\claim{c00.cl.hbar}{source-backed}
  \hbar = \frac{\hPlanck}{2\pi} .
\end{equation}
Die Zahlenwerte sind
$\hPlanck = \num[c00.cl.h-wert]{6{,}62607015 \cdot 10^{-34}}{J\,s}$ und
$\hbar = \num[c00.cl.hbar-wert]{1{,}054571817 \cdot 10^{-34}}{J\,s}$; der zweite ist auf
zehn Stellen abgebrochen.

Für $\nu > 0$ und $\Temp > 0$ betrachten wir den Exponenten
$x = \hPlanck\nu/(\kB \Temp)$. Für kleine $x$ ist
$e^x - 1$ nahezu gleich $x$. Setzt man das in \eqref{c00.eq.planck} ein, so kürzt sich
$\hPlanck$ heraus, und es bleibt die Funktion
\begin{equation}\label{c00.eq.grenzform}\claim{c00.cl.grenzform}{derived}
  u_{\mathrm{kl}}(\nu, \Temp) = \frac{8\pi \nu^2 \kB \Temp}{\clight^3},
  \qquad
  \frac{u(\nu, \Temp)}{u_{\mathrm{kl}}(\nu, \Temp)} = \frac{x}{e^{x} - 1}
  \quad \text{\tpcifenglish{with}{mit}}\ x = \frac{\hPlanck\nu}{\kB \Temp} .
\end{equation}
Die zweite Gleichung gilt exakt für $\nu > 0$ und $\Temp > 0$; man erhält sie, indem man
\eqref{c00.eq.planck} durch $u_{\mathrm{kl}}$ teilt. Das Verhältnis $x/(e^x - 1)$ liegt für
kleine $x$ nahe bei $1$ und fällt für große $x$ gegen $0$\zoomref{c00.v1}. Die Funktion
$u_{\mathrm{kl}}$ enthält $\hPlanck$ nicht mehr und ist hier die algebraisch gewonnene
Grenzform des im Skript angegebenen Strahlungsgesetzes
\srcref{tpc-script-v0.2}{vl01, Vorgeschichte, Punkt 1}.

\begin{tpcfigure}{c00.fig.planck}
  \caption{Das Strahlungsgesetz~\eqref{c00.eq.planck} und seine Grenzform
    \eqref{c00.eq.grenzform}. (a) Die Energiedichte $u$ für drei Temperaturen $T_0$,
    $\tfrac56 T_0$ und $\tfrac23 T_0$ (durchgezogen, gestrichelt, strichpunktiert),
    aufgetragen gegen $\hPlanck\nu/(\kB T_0)$ in Einheiten von
    $u_0 = 8\pi(\kB T_0)^3/(\hPlanck^2\clight^3)$. Die dünn gestrichelte Parabel ist
    die algebraische Grenzform $u_{\mathrm{kl}}$ für die höchste Temperatur $T_0$. Bei
    kleinen positiven Frequenzen liegt ihr Verhältnis zur durchgezogenen theoretischen
    Kurve nahe bei $1$; danach läuft die Parabel nach oben aus dem Bild.
    (b) Das Verhältnis $u/u_{\mathrm{kl}} = x/(e^x-1)$ für $\nu > 0$ und $\Temp > 0$.
    Es hängt nur von
    $x = \hPlanck\nu/(\kB\Temp)$ ab. Die schraffierte Fläche ist der Abstand zur
    Grenzform. Er wird merklich, sobald $\hPlanck\nu$ die Größe von $\kB\Temp$ erreicht.}
\end{tpcfigure}

Abbildung~\ref{c00.fig.planck} zeigt beides. Die Grenzform wächst wie $\nu^2$ ohne Ende,
während die Kurven des Strahlungsgesetzes ein Maximum haben und danach abfallen. Der
Unterschied wird sichtbar, sobald die Energie $\hPlanck\nu$ eines Lichtquants nicht mehr
klein gegen $\kB\Temp$ ist.

\begin{vertiefung}[id=c00.v1,
    titel={Die Grenzform, genau genommen},
    wozu={Der Kern sagt, das Verhältnis $x/(e^x-1)$ liege für kleine $x$ nahe bei $1$ und
      falle für große $x$ gegen $0$. Hier stehen zwei Schranken, die das beweisen, und
      eine Folgerung: Die Grenzform sagt eine unendliche Energiedichte voraus, das
      Strahlungsgesetz eine endliche.},
    voraussetzung={Kern bis Gleichung~\eqref{c00.eq.grenzform}; die Exponentialreihe und
      Integrale von Potenzen. Der Rest des Kapitels braucht diesen Block nicht.},
    zurueck={Für $\hPlanck\nu$ klein gegen $\kB\Temp$ stimmen Strahlungsgesetz und
      Grenzform überein; für große Frequenzen ist die Grenzform falsch, und zwar beliebig
      falsch.}]
  Für $x > 0$ ist $e^x = 1 + x + \tfrac12 x^2 + \tfrac16 x^3 + \dots$ eine Summe positiver
  Glieder. Daraus folgen zwei Abschätzungen für $e^x - 1$. Lässt man alle Glieder ab dem
  dritten weg, so wird die Summe kleiner: $e^x - 1 \geq x + \tfrac12 x^2$. Klammert man
  dagegen $x$ aus, so steht in der Klammer $1 + \tfrac12 x + \tfrac16 x^2 + \dots$, und
  jedes Glied davon ist höchstens so groß wie das entsprechende Glied
  $1 + x + \tfrac12 x^2 + \dots$ der Reihe für $e^x$: also $e^x - 1 \leq x\,e^x$.
  Zusammen ergibt das für alle $x > 0$
  \begin{equation}\label{c00.eq.schranken}\claim{c00.cl.schranken}{derived}
    e^{-x} \;\leq\; \frac{x}{e^{x} - 1} \;\leq\; \frac{1}{1 + \tfrac12 x} .
  \end{equation}
  Für $x \to 0$ gehen beide Schranken gegen $1$, für $x \to \infty$ beide gegen $0$. Bei
  $x = 1$, also $\hPlanck\nu = \kB\Temp$, liegt das Verhältnis zwischen $1/e$ und
  $\tfrac23$.

  Nun zur gesamten Energiedichte, dem Integral über alle Frequenzen. Für die Grenzform ist
  \begin{equation}\label{c00.eq.divergenz}\claim{c00.cl.divergenz}{derived}
    \int_0^{\nu_1} u_{\mathrm{kl}}(\nu, \Temp)\,\dd\nu
    = \frac{8\pi \kB \Temp}{3\,\clight^3}\,\nu_1^3 ,
  \end{equation}
  und das wächst mit der oberen Grenze $\nu_1$ über jede Schranke. Die Grenzform sagt für
  jede Temperatur $\Temp > 0$ eine unendliche Energiedichte voraus.

  Für das Strahlungsgesetz führt die Substitution $\nu = x\,\kB\Temp/\hPlanck$ auf
  \begin{equation}\label{c00.eq.gesamt}\claim{c00.cl.gesamt}{derived}
    \begin{aligned}
      \int_0^{\infty} u(\nu, \Temp)\,\dd\nu
      &= \frac{8\pi (\kB\Temp)^4}{\hPlanck^3 \clight^3} \int_0^\infty \frac{x^3}{e^x - 1}\,\dd x , \\
      \int_0^\infty \frac{x^3}{e^x - 1}\,\dd x &\leq \frac13 + 120 .
    \end{aligned}
  \end{equation}
  Die Schranke für das Integral erhält man in zwei Stücken. Für $0 < x \leq 1$ ist nach
  \eqref{c00.eq.schranken} $x^3/(e^x-1) \leq x^2$, und $\int_0^1 x^2\,\dd x = \tfrac13$.
  Für $x \geq 1$ benutzt man das Glied $x^5/120$ der Exponentialreihe:
  $e^x - 1 \geq x^5/120$, also $x^3/(e^x-1) \leq 120/x^2$, und
  $\int_1^\infty 120\,x^{-2}\,\dd x = 120$. Das Integral ist eine reine Zahl. Die gesamte
  Energiedichte des Strahlungsgesetzes ist also endlich und proportional zu $\Temp^4$.
  Den Wert der Zahl braucht dieser Kurs erst im Teil über statistische Physik.
\end{vertiefung}

\subsection{Photoeffekt}\label{c00.s4.2}

Licht kann nach dem Skript Elektronen auslösen. Für die Energie dieser Photoelektronen
gibt die Vorlesung die Beziehung an:
\begin{equation}\label{c00.eq.photoeffekt}\claim{c00.cl.photoeffekt}{source-backed}
  E_{\mathrm{el}} = \tfrac12 \massE v^2 = \hPlanck\nu - W .
\end{equation}
Dabei ist $\massE$ die Masse und $v$ die Geschwindigkeit des Elektrons, und $W$ ist seine
Bindungsenergie im Skript. Der Geltungsbereich dieser Beziehung wartet auf eine Quelle
oder Formulierung des Autors (G-0024). Die Vorlesung deutet die Beobachtung als Hinweis
auf Lichtquanten der Energie~\eqref{c00.eq.lichtquant}
\srcref{tpc-script-v0.2}{vl01, Vorgeschichte, Punkt 2}.

Eine kinetische Energie ist nie negativ. Gleichung~\eqref{c00.eq.photoeffekt} kann also
nur gelten, wenn
\begin{equation}\label{c00.eq.schwelle}\claim{c00.cl.schwelle}{derived}
  \hPlanck\nu \geq W,
  \qquad \text{\tpcifenglish{that is}{das heißt}} \qquad
  \nu \geq \frac{W}{\hPlanck} .
\end{equation}
Für festes $W$ beschreibt die angegebene Beziehung $E_{\mathrm{el}}$ als Gerade mit der
Steigung $\hPlanck$,
die bei $\nu = W/\hPlanck$ beginnt (Abbildung~\ref{c00.fig.photoeffekt}). Die Steigung ist
für jeden Wert von $W$ dieselbe.

\begin{tpcfigure}{c00.fig.photoeffekt}
  \caption{Die im Skript angegebene Beziehung~\eqref{c00.eq.photoeffekt}, algebraisch
    aufgetragen für zwei feste Werte der Bindungsenergie, $W_1$
    (durchgezogen) und $W_2 = \tfrac85 W_1$ (gestrichelt; das Verhältnis ist für das Bild
    gewählt). Beide Geraden haben dieselbe Steigung $\hPlanck$: Zu einem Frequenzunterschied
    $\nu' - \nu$ gehört der Energieunterschied $\hPlanck(\nu' - \nu)$. Jede Gerade beginnt bei
    ihrer Schwelle $\nu = W/\hPlanck$ aus Gleichung~\eqref{c00.eq.schwelle}. Die
    gepunkteten Verlängerungen treffen die Achse bei $-W$; dort ist die Energie negativ.
    Der physikalische Geltungsbereich der Beziehung steht aus (G-0024).}
\end{tpcfigure}

\subsection{Weitere Befunde}\label{c00.s4.3}

Die Vorlesung nennt vier weitere Befunde \srcref{tpc-script-v0.2}{vl01, Vorgeschichte, Punkte 3 bis 6}.
\begin{enumerate}
  \item \emph{Stabilität der Atome, Linienspektren, Franck-Hertz-Versuch.} Nach dem Skript
    wären Elektronenbahnen in der klassischen Physik instabil: Nach den
    Liénard-Wiechert-Potentialen der Elektrodynamik strahlt jede beschleunigte Punktladung.
    Ein Elektron auf einer Bahn um den Kern verlöre ständig Energie und müsste in den Kern
    stürzen, was es nicht tut. Linienspektren und die Versuche von Franck und Hertz deutet
    das Skript als Hinweis auf bestimmte, diskrete Energieniveaus in Atomen.
  \item \emph{Compton-Effekt.} Die Vorlesung behandelt die Streuung von Licht an
    Elektronen mit Lichtquanten der Energie~\eqref{c00.eq.lichtquant} und des Impulses $p$
    aus Gleichung~\eqref{c00.eq.impuls} und zieht Energie- und Impulserhaltung heran.
    Welchen Teil der Streuvorhersage dies erfasst, wartet auf eine Quelle oder
    Formulierung des Autors (G-0024).
  \item \emph{Stern-Gerlach-Versuch.} Das Skript beschreibt den Versuch mit endlich vielen
    Möglichkeiten der Ausrichtung des Drehimpulses relativ zu einem angelegten Magnetfeld.
    Der Geltungsbereich dieser Beschreibung wartet auf eine Quelle oder Formulierung
    des Autors (G-0024).
  \item \emph{Interferenz bei Elektronen.} Die Beugung von Elektronen an Kristallen
    (Davisson-Germer-Experiment) ist ein Hinweis auf Welleneigenschaften des Elektrons.
\end{enumerate}
Der Impuls eines Lichtquants der Wellenlänge $\lambda$ ist dabei
\begin{equation}\label{c00.eq.impuls}\claim{c00.cl.impuls}{source-backed}
  p = \frac{\hPlanck}{\lambda} .
\end{equation}
Das Linienspektrum des Wasserstoffs wird der Kurs in Kapitel~11 ausführlich behandeln.
Jahreszahlen und Namen zu diesen Befunden gehören nicht zum Kern\zoomref{c00.v2}.

\begin{vertiefung}[id=c00.v2,
    titel={Namen und Jahreszahlen, wie die Vorlesung sie nennt},
    wozu={Der Kern nennt die Befunde ohne Geschichte. Hier stehen die Jahreszahlen und
      Namen aus dem Vorlesungsskript, für alle, die die Geschichte erzählen wollen.},
    voraussetzung={Abschnitt~\ref{c00.s4}. Der Rest des Kurses braucht diesen Block
      nicht.},
    zurueck={Die mathematische Form der Quantenmechanik und ihre statistische Deutung
      stammen nach dem Skript aus den Jahren 1925 bis 1927.}]
  Die Angaben dieses Blocks sind dem Vorlesungsskript entnommen und nicht gegen eine
  historische Quelle geprüft \srcref{tpc-script-v0.2}{vl01, Vorgeschichte}. Sein
  Ausgangspunkt:
  \begin{equation}\label{c00.eq.geburtsjahr}\claim{c00.cl.geschichte}{source-backed}
    \text{\tpcifenglish{The year of birth of quantum mechanics is 1926.}{Das Geburtsjahr der Quantenmechanik ist 1926.}}
  \end{equation}

  \begin{tpctable}{c00.tab.jahre}
    \caption{Jahreszahlen und Namen zu den Befunden des Abschnitts~\ref{c00.s4}, wie das
      Vorlesungsskript sie angibt.}
    \begin{tabular}{ll}
      \toprule
      Befund & Angabe des Skripts \\
      \midrule
      Strahlungsgesetz, Konstante $\hPlanck$ & Planck, 1899 \\
      Photoeffekt & Lenard, 1899; Einstein, 1905 \\
      Franck-Hertz-Versuch & um 1911 \\
      Stern-Gerlach-Versuch & 1922 \\
      Welleneigenschaften des Elektrons & de Broglie, 1924 (Postulat) \\
      Davisson-Germer-Experiment & 1929 \\
      \bottomrule
    \end{tabular}
  \end{tpctable}

  Planck führte die Konstante $\hPlanck$ nach dem Skript ursprünglich als „Hilfsgröße“ ein.
  Im Jahr 1926 veröffentlichte Schrödinger eine Reihe von Arbeiten über Wellenmechanik.
  Ein paralleler Zugang begann 1925 mit Überlegungen von Heisenberg; daraus entstanden die
  Matrizenmechanik von Born, Jordan und Heisenberg und eine Fassung von Dirac. Noch 1926
  zeigte Schrödinger die Äquivalenz beider Zugänge, und Born lieferte die statistische
  Interpretation der Theorie. Der Autor der Vorlesung hält Borns Beitrag für den
  wichtigsten Aspekt der ganzen Theorie.

  Fast gleichzeitig erhielt die Theorie ihre heutige mathematische Form. Born ließ sich
  von Norbert Wiener beraten, wie man aus Matrizen Operatoren macht. Hilbert hielt im
  Wintersemester 1926/27 in Göttingen eine Vorlesung über Quantenmechanik, und sein
  Assistent John von Neumann führte den Hilbertraum, Dichtematrizen und die Entropie für
  Quantenzustände ein. Die Rolle der Gruppentheorie klärten Wigner und Weyl. Die
  Interpretation der Theorie blieb umstritten und ist es noch; einen Widerspruch zur
  Quantenmechanik gibt es nach dem Skript bis heute nicht.

  \begin{historical}
    Wer Deutsch liest, hat Glück: Die Originalarbeiten von Planck, Schrödinger, Heisenberg
    und Born sind auf Deutsch geschrieben.
  \end{historical}
\end{vertiefung}

\subsection{Zwei Sätze der Vorlesung}\label{c00.s4.4}

Die Vorlesung fasst die Befunde in zwei Sätzen zusammen
\srcref{tpc-script-v0.2}{vl01, die beiden gelben Kästen}. Der erste:
\begin{equation}\label{c00.eq.welle-teilchen}\claim{c00.cl.welle-teilchen}{source-quoted}
  \begin{aligned}
    &\text{\tpcifenglish{Matter and radiation possess}{Materie und Strahlung besitzen}} \\
    &\text{\tpcifenglish{both wave and particle properties.}{sowohl Wellen- als auch Teilcheneigenschaften.}}
  \end{aligned}
\end{equation}
Der zweite:
\begin{equation}\label{c00.eq.kernsatz}\claim{c00.cl.kernsatz}{source-quoted}
  \begin{aligned}
    &\text{\tpcifenglish{Quantum mechanics is a probabilistic theory.}{Die Quantenmechanik ist eine probabilistische Theorie.}} \\
    &\text{\tpcifenglish{It therefore gives us predictions}{Sie liefert uns somit nur Vorhersagen}} \\
    &\text{\tpcifenglish{only in the form of probabilities.}{in Form von Wahrscheinlichkeiten.}}
  \end{aligned}
\end{equation}
Der zweite Satz ist dem ersten übergeordnet. „Wellen- und Teilcheneigenschaften“ bleibt ein
wenig schwammig; die probabilistische Aussage ist präzise und liegt allen Quantenphänomenen
zugrunde. In der Vorlesung heißt es dazu: In jeder mündlichen Prüfung wird gefragt, was die
wichtigste Aussage über die Quantenmechanik ist.

\begin{merksatz}[Merksatz]
  Die Quantenmechanik ist eine probabilistische Theorie: Sie sagt Wahrscheinlichkeiten
  voraus. Das ist Satz~\eqref{c00.eq.kernsatz} der Vorlesung. Der Kapitelautor schlägt
  vor, ihn mit P0 in Satz~\eqref{c00.eq.p0} für die Quantenmechanik zu verbinden;
  diese Lesart wartet auf die Entscheidung des Autors (G-0023).
\end{merksatz}

Das Skript stellt dem die klassische Physik gegenüber: Dort lauten die Vorhersagen
definitiv, das Teilchen ist zu dieser Zeit an dieser Stelle. Der Kapitelautor schlägt
als endliches Beispiel zu Satz~\eqref{c00.eq.probabilistisch} die Punktgewichte aus
Kapitel~1 vor: Eine solche Präparation legt ihr ganzes Gewicht auf ein einziges Ergebnis.
Diese Deutung wartet auf die Entscheidung des Autors (G-0023); der Satz über
Präparationen legt keine Dynamik deterministischer Theorien fest.

Die Zahlen, die die Quantenmechanik liefert, sind Wahrscheinlichkeiten. Bevor ihre
mathematischen Grundlagen aufgebaut werden können, muss deshalb geklärt sein, wie man mit
Wahrscheinlichkeiten rechnet. Das tut das nächste Kapitel.
